\documentclass{article}
\usepackage{graphicx}
\usepackage{booktabs}
\usepackage{amsmath}
\usepackage{hyperref}
\usepackage{multirow}

\title{EnterpriseSynth: Generating Verified Agentic Training Data from Proprietary API Schemas Without Live Execution}
\author{[Intern Name], Spurthi Setty \\
Anote AI}
\date{}

\begin{document}

\maketitle

\begin{abstract}
\noindent
Fine-tuning small language models for enterprise tool-calling requires high-quality agentic traces, but enterprises cannot use public API datasets (privacy, domain mismatch) and cannot batch-execute live internal APIs (side effects, cost, auth). We present \textbf{EnterpriseSynth}, a framework that accepts any OpenAPI/Swagger specification and produces two verified outputs: (1) supervised fine-tuning traces and (2) evaluation records with intent specifications --- without executing a single live API call. We compare three generation strategies: schema-constrained decoding, multi-step mock-execution verification, and adversarial self-play. We validate on three public enterprise-adjacent API schemas (CRM, Financial, DevOps), fine-tune Qwen 2.5 7B and Llama 3 8B on each strategy's output, and show that synthetic SFT improves tool-call performance by X\% over zero-shot. EnterpriseSynth directly produces the data format required by EnterpriseBench \cite{agenticeval}.
\end{abstract}

\section{Introduction}
% The enterprise cold-start problem: no training data for proprietary workflows
% Prior work (APIGen, TOUCAN) requires live/public APIs
% Our solution: schema-only, dual-output (SFT + eval), enterprise-safe

\section{Related Work}
\subsection{Agentic Data Generation}
% APIGen (NeurIPS 2024), APIGen-MT (NeurIPS 2025), TOUCAN (2025), ToolLLM
\subsection{Synthetic Data for Fine-Tuning}
\subsection{Schema-Constrained Generation}
% LMQL, Outlines, Guidance

\begin{table}[h!]
\centering
\caption{Comparison with Prior Work}
\begin{tabular}{lccccc}
\toprule
\textbf{Paper} & \textbf{Schema-Only} & \textbf{Multi-Step} & \textbf{Dual Output} & \textbf{Intent Spec} & \textbf{Enterprise} \\
\midrule
ToolLLM (ICLR 2024) & \texttimes & \texttimes & \texttimes & \texttimes & \texttimes \\
APIGen (NeurIPS 2024) & \texttimes & \texttimes & \texttimes & \texttimes & \texttimes \\
APIGen-MT (NeurIPS 2025) & \texttimes & \checkmark & \texttimes & $\sim$ & \texttimes \\
TOUCAN (2025) & \texttimes & \checkmark & \texttimes & \texttimes & \texttimes \\
\textbf{EnterpriseSynth (ours)} & \checkmark & \checkmark & \checkmark & \checkmark & \checkmark \\
\bottomrule
\end{tabular}
\end{table}

\section{EnterpriseSynth Framework}
\subsection{Input: OpenAPI/Swagger Schema Parsing}
\subsection{Generation Strategies}
% Strategy A: Schema-Constrained Decoding
% Strategy B: Multi-Step Verification Pipeline (mock execution)
% Strategy C: Self-Play Adversarial Multi-Agent
% Baseline: Pure LLM generation (no constraints)
\subsection{Dual Output Architecture}
% SFT Record format: {instruction, tool_schema, verified_trace}
% Eval Record format: {instruction, intent_specification, ground_truth_trace, scope_constraints}

\section{Experimental Setup}
\subsection{API Schemas}
% HubSpot (CRM), Plaid-adjacent (Finance), GitHub Actions (DevOps)
\subsection{Fine-Tuning Protocol}
% Qwen 2.5 7B, Llama 3 8B
\subsection{Evaluation}
% Zero-shot → SFT performance gap on EnterpriseBench tasks

\section{Results}
% Table: Strategy comparison (syntactic validity, semantic validity, diversity, cost)
% Table: Fine-tuned model performance on EnterpriseBench
% Figure: Performance gap: zero-shot vs. each strategy's SFT output

\section{Analysis}
\subsection{Quality of Generated Traces}
\subsection{Mock Execution Fidelity}
\subsection{Self-Play Adversarial Pressure Analysis}

\section{Conclusion}

\section*{References}
\small
\begin{enumerate}
    \item Liu et al. APIGen: Automated Pipeline for Generating Verifiable and Diverse Function-Calling Datasets, 2024. NeurIPS 2024 D\&B.
    \item Prabhakar et al. APIGen-MT, 2025. NeurIPS 2025 D\&B.
    \item Xu et al. TOUCAN, arXiv Oct 2025.
    % Add more references
\end{enumerate}

\end{document}
